\documentclass[11pt,a4paper]{article}
\usepackage[top=2.3cm,bottom=2.3cm,left=2.2cm,right=2.2cm]{geometry}

\usepackage{amsmath,amssymb}
\usepackage{fontspec}
\setmainfont{Liberation Serif}
\setsansfont{Liberation Sans}
\setmonofont{DejaVu Sans Mono}[Scale=0.84]

\usepackage{xcolor}
\definecolor{brandblue}{RGB}{20, 55, 105}
\definecolor{accentblue}{RGB}{35, 95, 165}
\definecolor{codebg}{RGB}{248, 249, 251}
\definecolor{codeframe}{RGB}{215, 222, 230}
\definecolor{javakeyword}{RGB}{0, 45, 170}
\definecolor{javastring}{RGB}{5, 120, 20}
\definecolor{javacomment}{RGB}{120, 120, 120}
\definecolor{javanumber}{RGB}{20, 75, 220}

\definecolor{noteborder}{RGB}{30, 110, 65}
\definecolor{notebg}{RGB}{242, 249, 245}
\definecolor{warnborder}{RGB}{185, 75, 10}
\definecolor{warnbg}{RGB}{254, 248, 240}
\definecolor{tipborder}{RGB}{30, 95, 160}
\definecolor{tipbg}{RGB}{242, 247, 254}

\usepackage{fancyhdr}
\usepackage{lastpage}
\pagestyle{fancy}
\fancyhf{}
\lhead{\parbox[t]{0.58\textwidth}{\raggedright\small\sffamily\color{brandblue}\textbf{T19 --- Documentazione con Javadoc e Unit Testing con JUnit 5}}}
\rhead{\small\sffamily\color{gray}Programmazione II -- UniVR (2025/2026)}
\cfoot{\small\sffamily Pagina \thepage\ di \pageref{LastPage}}
\renewcommand{\headrulewidth}{0.5pt}
\renewcommand{\headrule}{\hbox to\headwidth{\color{brandblue!70!black}\leaders\hrule height \headrulewidth\hfill}}

\usepackage{titlesec}
\titleformat{\section}{\Large\bfseries\sffamily\color{brandblue}}{\thesection}{0.8em}{}[\titlerule]
\titleformat{\subsection}{\large\bfseries\sffamily\color{accentblue}}{\thesubsection}{0.8em}{}
\titleformat{\subsubsection}{\normalsize\bfseries\sffamily\color{brandblue!85!black}}{\thesubsubsection}{0.8em}{}

\usepackage{needspace}
\usepackage{fancyvrb}
\DefineVerbatimEnvironment{asciibox}{Verbatim}{
  fontsize=\fontsize{7.5pt}{9.2pt}\selectfont,
  frame=single,
  rulecolor=\color{codeframe},
  fillcolor=\color{codebg},
  framesep=2.5mm
}
\usepackage{listings}
\lstset{
  backgroundcolor=\color{codebg},
  frame=single,
  rulecolor=\color{codeframe},
  basicstyle=\ttfamily\footnotesize,
  keywordstyle=\color{javakeyword}\bfseries,
  stringstyle=\color{javastring},
  commentstyle=\color{javacomment}\itshape,
  numberstyle=\tiny\color{gray},
  numbers=left,
  numbersep=7pt,
  stepnumber=1,
  breaklines=true,
  breakatwhitespace=true,
  showstringspaces=false,
  tabsize=4,
  xleftmargin=12pt,
  framexleftmargin=12pt
}

\newcommand{\passthrough}[1]{#1}
\providecommand{\tightlist}{\setlength{\itemsep}{0pt}\setlength{\parskip}{0pt}}

\usepackage{tcolorbox}
\tcbuselibrary{skins,breakable}
\newtcolorbox{studynote}[1][Nota]{
  enhanced,
  breakable,
  colback=notebg,
  colframe=noteborder,
  fonttitle=\bfseries\sffamily\small,
  coltitle=white,
  title={#1},
  arc=2mm,
  boxrule=0.8pt,
  left=9pt, right=9pt, top=7pt, bottom=7pt
}

\newtcolorbox{studywarning}[1][Attenzione]{
  enhanced,
  breakable,
  colback=warnbg,
  colframe=warnborder,
  fonttitle=\bfseries\sffamily\small,
  coltitle=white,
  title={#1},
  arc=2mm,
  boxrule=0.8pt,
  left=9pt, right=9pt, top=7pt, bottom=7pt
}

\newtcolorbox{studytip}[1][Suggerimento]{
  enhanced,
  breakable,
  colback=tipbg,
  colframe=tipborder,
  fonttitle=\bfseries\sffamily\small,
  coltitle=white,
  title={#1},
  arc=2mm,
  boxrule=0.8pt,
  left=9pt, right=9pt, top=7pt, bottom=7pt
}

\usepackage{booktabs}
\usepackage{longtable}
\usepackage{array}
\usepackage{calc}

\usepackage{hyperref}
\hypersetup{
  colorlinks=true,
  linkcolor=accentblue,
  urlcolor=accentblue,
  citecolor=brandblue,
  pdfborder={0 0 0}
}

\title{\Huge\bfseries\sffamily\color{brandblue} T19 --- Documentazione con Javadoc e Unit Testing con JUnit 5 \\[0.3em] \large\sffamily Convenzioni Javadoc, Architettura JUnit 5, Annotazioni di Ciclo di Vita ed Asserzioni}
\author{\textbf{Docente:} Prof. Michele Pasqua \quad---\quad \textbf{Appunti Revisione:} Wally}
\date{A.A. 2025/2026 -- Universit\`a degli Studi di Verona}

\begin{document}
\maketitle
\thispagestyle{fancy}
\vspace{-1.2em}
\noindent\textcolor{brandblue}{\rule{\textwidth}{0.8pt}}
\vspace{1.2em}

\begin{center}\rule{0.5\linewidth}{0.5pt}\end{center}

\subsection{1. Analisi Teorica Approfondita (Slide
T19)}\label{analisi-teorica-approfondita-slide-t19}

La lezione 19 approfondisce due pratiche ingegneristiche irrinunciabili
nello sviluppo professionale in Java: 1. \textbf{Documentazione
automatizzata del codice} tramite lo standard \textbf{Javadoc}. 2.
\textbf{Collaudo e Unit Testing automatizzato} tramite il framework
\textbf{JUnit 5 (Jupiter)} e il plugin Maven \textbf{Surefire}.

\begin{center}\rule{0.5\linewidth}{0.5pt}\end{center}

\subsubsection{1.1 Documentazione con Javadoc (Slide
3-11)}\label{documentazione-con-javadoc-slide-3-11}

La documentazione del codice sorgente è essenziale per manutenibilità,
collaborazione e pubblicazione di API riusabili. In Java la
documentazione vive a stretto contatto con il codice nei cosiddetti
\textbf{Doc Comments}: * \textbf{Sintassi}: Delimitati da
\passthrough{\lstinline!/**!} all'inizio e \passthrough{\lstinline!*/!}
alla fine. Ogni riga intermedia inizia convenzionalmente con
\passthrough{\lstinline!*!}:
\passthrough{\lstinline!java   /**    * Descrizione sintetica del componente.    * <p>Supporta l'uso di tag HTML come paragrafi, link e formattazione.</p>    */!}
* \textbf{I Tag Javadoc Standard (Slide 5-6)}: Iniziano con
\passthrough{\lstinline!@!}, sono \emph{case-sensitive} e devono
comparire a inizio riga: * \passthrough{\lstinline!@param <nome>!}:
descrive un parametro formale di un metodo o costruttore. *
\passthrough{\lstinline!@return!}: descrive il significato del valore
restituito da un metodo (omesso nei metodi
\passthrough{\lstinline!void!} e nei costruttori). *
\passthrough{\lstinline!@throws <Eccezione>!} (o
\passthrough{\lstinline!@exception!}): descrive le condizioni anomale
che causano il sollevamento di una determinata eccezione. *
\passthrough{\lstinline!\{@link <package.Classe\#metodo>\}!}: genera un
ipertesto cliccabile verso un'altra classe o metodo. *
\passthrough{\lstinline!@author!}: autore del modulo. *
\passthrough{\lstinline!@since <versione>!}: versione a partire dalla
quale la feature è disponibile. * \passthrough{\lstinline!@version!}:
versione corrente del sorgente. * \passthrough{\lstinline!@deprecated!}:
avvisa che l'elemento è obsoleto, illustrandone il motivo e
l'alternativa moderna da utilizzare.

\paragraph{Regole di Visibilità e Generazione (Slide
9-11)}\label{regole-di-visibilituxe0-e-generazione-slide-9-11}

\begin{itemize}
\item
  \textbf{Default di Javadoc}: Per impostazione predefinita, Javadoc
  genera la documentazione per le sole entità
  \textbf{\passthrough{\lstinline!public!} e
  \passthrough{\lstinline!protected!}} (l'interfaccia pubblica esposta
  ai client). I campi e metodi \passthrough{\lstinline!private!} o
  \emph{package-private} vengono omessi.
\item
  \textbf{Inclusione del privato}: Si deve passare esplicitamente il
  flag \passthrough{\lstinline!-private!} da riga di comando:

\begin{lstlisting}[language=bash]
javadoc -private -d doc *
\end{lstlisting}
\item
  \textbf{Integrazione Maven
  (\passthrough{\lstinline!maven-javadoc-plugin!}, Slide 11)}:
  Aggiungendo il plugin nel \passthrough{\lstinline!pom.xml!}, la
  documentazione viene generata con un singolo comando standard:

\begin{lstlisting}[language=bash]
mvn javadoc:javadoc
\end{lstlisting}
\end{itemize}

\begin{center}\rule{0.5\linewidth}{0.5pt}\end{center}

\subsubsection{1.2 Unit Testing con JUnit (Slide
13-16)}\label{unit-testing-con-junit-slide-13-16}

Il testing programmatico garantisce la correttezza delle singole unità
software in isolamento (metodi o singole classi): * \textbf{Vantaggi
sistemici}: * \textbf{Test-Driven Development (TDD)}: scrivere i test
prima ancora di implementare il codice applicativo. * \textbf{Regression
Testing}: certezza che refactoring o nuove feature non rompano
comportamenti pregressi funzionanti. * Riduzione drastica del tempo
speso nel debugging manuale. * \textbf{Anatomia di un Test Case}: 1.
Preparazione dell'input (\emph{Arrange / Setup}). 2. Esecuzione del
metodo sotto test (\emph{Act}). 3. Confronto dell'output effettivo
(\emph{Actual}) con l'output atteso (\emph{Expected}) tramite
\textbf{asserzioni} (\emph{Assert}). * \textbf{Meccanismo di
Fallimento}: * Un metodo di test restituisce sempre
\passthrough{\lstinline!void!}. * Se tutte le asserzioni sono
soddisfatte, il metodo termina normalmente e JUnit lo contrassegna come
\textbf{PASSED} (verde). * Se un'asserzione fallisce, solleva un errore
di tipo \passthrough{\lstinline!AssertionError!} (o
\passthrough{\lstinline!AssertionFailedError!}). Il framework JUnit
cattura l'errore, contrassegna il test come \textbf{FAILED} (rosso) con
il report della discrepanza, e \textbf{prosegue regolarmente con i test
successivi} senza arrestare l'intera suite.

\begin{center}\rule{0.5\linewidth}{0.5pt}\end{center}

\subsubsection{\texorpdfstring{1.3 Asserzioni in JUnit 5
(\texttt{org.junit.jupiter.api.Assertions}, Slide
17-21)}{1.3 Asserzioni in JUnit 5 (org.junit.jupiter.api.Assertions, Slide 17-21)}}\label{asserzioni-in-junit-5-org.junit.jupiter.api.assertions-slide-17-21}

Il pacchetto Jupiter standardizza una vasta famiglia di metodi statici:

\begin{longtable}[]{@{}
  >{\raggedright\arraybackslash}p{(\linewidth - 2\tabcolsep) * \real{0.5000}}
  >{\raggedright\arraybackslash}p{(\linewidth - 2\tabcolsep) * \real{0.5000}}@{}}
\toprule\noalign{}
\begin{minipage}[b]{\linewidth}\raggedright
Metodo Asserzione
\end{minipage} & \begin{minipage}[b]{\linewidth}\raggedright
Comportamento
\end{minipage} \\
\midrule\noalign{}
\endhead
\bottomrule\noalign{}
\endlastfoot
\passthrough{\lstinline!assertEquals(expected, actual, [msg])!} &
Verifica che \passthrough{\lstinline!expected.equals(actual)!}. Valido
per primitivi e oggetti. \\
\passthrough{\lstinline!assertTrue(condition, [msg])!} & Verifica che la
condizione booleana sia \passthrough{\lstinline!true!}. \\
\passthrough{\lstinline!assertFalse(condition, [msg])!} & Verifica che
la condizione booleana sia \passthrough{\lstinline!false!}. \\
\passthrough{\lstinline!assertArrayEquals(expected, actual)!} &
Confronta il contenuto e l'ordinamento di due array elemento per
elemento. \\
\passthrough{\lstinline!assertNull(obj)!} /
\passthrough{\lstinline!assertNotNull(obj)!} & Verifica se il puntatore
è \passthrough{\lstinline!null!} o non nullo. \\
\passthrough{\lstinline!assertThrows(Exception.class, executable)!} &
Verifica che l'esecuzione di una lambda/blocco sollevi tassativamente
l'eccezione attesa! \\
\end{longtable}

\paragraph{Best Practice: Static Imports (Slide
21)}\label{best-practice-static-imports-slide-21}

Per evitare di anteporre continuamente
\passthrough{\lstinline!Assertions.!} davanti a ogni chiamata, si usa
l'import statico:

\begin{lstlisting}[language=Java]
import static org.junit.jupiter.api.Assertions.assertEquals;
import static org.junit.jupiter.api.Assertions.assertThrows;
\end{lstlisting}

\begin{center}\rule{0.5\linewidth}{0.5pt}\end{center}

\subsubsection{1.4 Il Ciclo di Vita del Test e le Annotazioni JUnit 5
(Slide
22)}\label{il-ciclo-di-vita-del-test-e-le-annotazioni-junit-5-slide-22}

\needspace{7cm}
\begin{asciibox}
┌────────────────────────┐
│      @BeforeAll        │ (eseguito 1 sola volta prima di tutti, STATIC)
└───────────┬────────────┘
            │
┌───────────┴─────────────────┐
│   Per ogni metodo @Test:    │
│  ┌────────────────────────┐ │
│  │      @BeforeEach       │ │ (setup specifico per il test)
│  └───────────┬────────────┘ │
│              ▼              │
│  ┌────────────────────────┐ │
│  │         @Test          │ │ (esecuzione del caso di test)
│  └───────────┬────────────┘ │
│              ▼              │
│  ┌────────────────────────┐ │
│  │       @AfterEach       │ │ (teardown / pulizia risorse)
│  └────────────────────────┘ │
└───────────┬─────────────────┘
            │
┌───────────┴────────────┐
│       @AfterAll        │ (eseguito 1 sola volta alla fine, STATIC)
└────────────────────────┘
\end{asciibox}

\begin{itemize}
\tightlist
\item
  \passthrough{\lstinline!@BeforeAll!}: Inizializzazione ``pesante'' o
  condivisa tra tutti i test (es. apertura connessione, setup di un
  database o di costanti immutabili). \textbf{Deve essere
  \passthrough{\lstinline!static!}}.
\item
  \passthrough{\lstinline!@AfterAll!}: Chiusura finale delle risorse
  globali. \textbf{Deve essere \passthrough{\lstinline!static!}}.
\item
  \passthrough{\lstinline!@BeforeEach!}: Inizializzazione fresca prima
  di \emph{ciascun} test per garantire l'isolamento e l'indipendenza dei
  test (evitando che un test modifichi lo stato di un altro).
\item
  \passthrough{\lstinline!@AfterEach!}: Pulizia post-test (es.
  cancellazione di file temporanei creati durante il test).
\item
  \passthrough{\lstinline!@DisplayName("Descrizione chiara")!}:
  Personalizza il nome del test visualizzato nei report o nell'IDE.
\item
  \passthrough{\lstinline!@Disabled!}: Disabilita temporaneamente il
  test senza doverlo commentare o cancellare.
\end{itemize}

\begin{center}\rule{0.5\linewidth}{0.5pt}\end{center}

\subsubsection{1.5 Organizzazione del Progetto e Convenzioni Maven
(Slide
23-24)}\label{organizzazione-del-progetto-e-convenzioni-maven-slide-23-24}

Maven e i moderni build tool impongono una struttura standard di
cartelle: * \textbf{\passthrough{\lstinline!src/main/java!}}: Contiene
il codice sorgente dell'applicazione (es.
\passthrough{\lstinline!it.oop.core.Date!}). *
\textbf{\passthrough{\lstinline!src/test/java!}}: Contiene
esclusivamente le classi di collaudo. * \textbf{Allineamento dei Package
(Regola Fondamentale)}: Una classe di test che verifica
\passthrough{\lstinline!it.oop.core.Date!} \textbf{deve risiedere nello
stesso identico package \passthrough{\lstinline!it.oop.core!}}
(all'interno di \passthrough{\lstinline!src/test/java!}). In questo modo
la classe di test ha accesso non solo ai membri
\passthrough{\lstinline!public!}, ma anche a tutti i metodi e campi con
visibilità di \textbf{package
(\passthrough{\lstinline!package-private!})}, facilitando il testing
interno senza forzare l'apertura a \passthrough{\lstinline!public!} di
dettagli architetturali riservati!

\begin{center}\rule{0.5\linewidth}{0.5pt}\end{center}

\subsection{\texorpdfstring{2. Analisi Dettagliata del Codice
(\texttt{Lezione18/MavenDate})}{2. Analisi Dettagliata del Codice (Lezione18/MavenDate)}}\label{analisi-dettagliata-del-codice-lezione18mavendate}

In \passthrough{\lstinline!Lezione18/MavenDate!} troviamo la
configurazione completa del plugin Maven Surefire e due classi di test
che collaudano la gerarchia di \passthrough{\lstinline!Date!}.

\subsubsection{\texorpdfstring{2.1 Configurazione in
\href{file:///home/wally/Documenti/GitHub/Programmazione-II/25-26/Teoria-e-Esercitazioni/Codice-20260902/Lezione18/MavenDate/pom.xml\#L44-L68}{pom.xml}}{2.1 Configurazione in pom.xml}}\label{configurazione-in-pom.xml}

\begin{lstlisting}[language=XML]
    <build>
        <plugins>
            <!-- Plugin Javadoc (Slide 11) -->
            <plugin>
                <groupId>org.apache.maven.plugins</groupId>
                <artifactId>maven-javadoc-plugin</artifactId>
                <version>3.6.2</version>
                <configuration>
                    <source>1.8</source>
                    <show>private</show> <!-- Include metodi e campi privati -->
                </configuration>
            </plugin>
        </plugins>
    </build>

    <dependencies>
        <!-- Motore di esecuzione JUnit 5 (Slide 24) -->
        <dependency>
            <groupId>org.junit.jupiter</groupId>
            <artifactId>junit-jupiter-engine</artifactId>
            <version>5.10.0</version>
            <scope>test</scope> <!-- Visibile solo durante la fase di test -->
        </dependency>
        <!-- Maven Surefire Plugin per l'esecuzione di 'mvn test' -->
        <dependency>
            <groupId>org.apache.maven.plugins</groupId>
            <artifactId>maven-surefire-plugin</artifactId>
            <version>3.5.4</version>
        </dependency>
    </dependencies>
\end{lstlisting}

\begin{center}\rule{0.5\linewidth}{0.5pt}\end{center}

\subsubsection{\texorpdfstring{2.2
\href{file:///home/wally/Documenti/GitHub/Programmazione-II/25-26/Teoria-e-Esercitazioni/Codice-20260902/Lezione18/MavenDate/src/test/java/it/oop/core/TestDate.java}{TestDate.java}}{2.2 TestDate.java}}\label{testdate.java}

Questa classe collauda sia il flusso nominale (costruzione valida) che
il flusso eccezionale (lancio di eccezioni):

\begin{lstlisting}[language=Java]
package it.oop.core;

import it.oop.exception.IllegalDateException;
import org.junit.jupiter.api.Assertions;
import org.junit.jupiter.api.Test;
import static org.junit.jupiter.api.Assertions.assertThrows;

public class TestDate {

    private Date date;

    // 1. Collaudo del caso nominale con assertEquals
    @Test
    public void testDateConstructor() {
        date = new Date(1, 12, 2025);
        Assertions.assertEquals(1, date.getDay());
        Assertions.assertEquals(12, date.getMonth());
        Assertions.assertEquals(2025, date.getYear());
    }

    // 2. Collaudo del caso eccezionale con assertThrows e Lambda
    @Test
    public void testDateConstructorException() {
        // Verifica che passando un anno negativo (-2), il costruttore sollevi tassativamente IllegalDateException
        assertThrows(IllegalDateException.class, () -> date = new Date(1, 12, -2));
    }
}
\end{lstlisting}

\begin{center}\rule{0.5\linewidth}{0.5pt}\end{center}

\subsubsection{\texorpdfstring{2.3
\href{file:///home/wally/Documenti/GitHub/Programmazione-II/25-26/Teoria-e-Esercitazioni/Codice-20260902/Lezione18/MavenDate/src/test/java/it/oop/core/TestItalianDate.java}{TestItalianDate.java}}{2.3 TestItalianDate.java}}\label{testitaliandate.java}

Esemplifica l'uso della fixture \passthrough{\lstinline!@BeforeAll!} e
mette in luce il superamento della vecchia parola chiave
\passthrough{\lstinline!assert!}:

\begin{lstlisting}[language=Java]
package it.oop.core;

import org.junit.jupiter.api.Assertions;
import org.junit.jupiter.api.BeforeAll;
import org.junit.jupiter.api.Test;

public class TestItalianDate {
    private static ItalianDate date;

    // Fixture statica eseguita una volta sola prima di tutti i test
    @BeforeAll
    public static void setup() {
        date = new ItalianDate(1, 1, 1970);
    }

    @Test
    public void printFormatTest() {
        // Verifica del formato con JUnit 5:
        Assertions.assertEquals("dd/mm/yyyy", date.printFormat());
        
        // Confronto con la vecchia sintassi 'assert' del linguaggio (commentata dal docente):
        // assert date.printFormat().equals("dd/mm/yyyy") : "wrong format";
    }
}
\end{lstlisting}

\begin{center}\rule{0.5\linewidth}{0.5pt}\end{center}

\subsection{3. Risultato di Esecuzione Reale da
Terminale}\label{risultato-di-esecuzione-reale-da-terminale}

Comando eseguito nel progetto
\href{file:///home/wally/Documenti/GitHub/Programmazione-II/25-26/Teoria-e-Esercitazioni/Codice-20260902/Lezione18/MavenDate}{Lezione18/MavenDate}:

\begin{lstlisting}[language=bash]
mvn test
\end{lstlisting}

Output ottenuto:

\begin{lstlisting}
[INFO] -------------------------------------------------------
[INFO]  T E S T S
[INFO] -------------------------------------------------------
[INFO] Running it.oop.core.TestDate
[INFO] Tests run: 2, Failures: 0, Errors: 0, Skipped: 0, Time elapsed: 0.106 s -- in it.oop.core.TestDate
[INFO] Running it.oop.core.TestItalianDate
[INFO] Tests run: 1, Failures: 0, Errors: 0, Skipped: 0, Time elapsed: 0.021 s -- in it.oop.core.TestItalianDate
[INFO] 
[INFO] Results:
[INFO] 
[INFO] Tests run: 3, Failures: 0, Errors: 0, Skipped: 0
[INFO] 
[INFO] ------------------------------------------------------------------------
[INFO] BUILD SUCCESS
[INFO] ------------------------------------------------------------------------
\end{lstlisting}

I 3 test sono stati eseguiti con successo, validando in modo formale la
suite del progetto e generando i relativi file di report XML e TXT nella
cartella \passthrough{\lstinline!target/surefire-reports/!}.

\begin{center}\rule{0.5\linewidth}{0.5pt}\end{center}

Dimmi \textbf{``vai''} per procedere con il \textbf{Blocco 19} (Lezione
19 / Slide T20 --- \emph{Java I/O}, stream di byte vs caratteri,
\passthrough{\lstinline!InputStream!},
\passthrough{\lstinline!OutputStream!},
\passthrough{\lstinline!Reader!}, \passthrough{\lstinline!Writer!},
\passthrough{\lstinline!BufferedReader!}, serializzazione, gestione
delle risorse con \passthrough{\lstinline!try-with-resources!} e analisi
di \passthrough{\lstinline!MainIO.java!} in
\passthrough{\lstinline!Lezione19/MavenDate!}).


\end{document}
